\section{Nutritional adequacy and yield-to-diet methods}\label{sec:nutmethod}

\subsection{Nutrient adequacy}\label{sec:adequacy}
For country $c$ in year $t$, national supply $S_n$ (per person per day) is compared with a requirement $R_n$:
\begin{equation}\label{eq:nar}
\text{NAR}_n = \min\!\left(1,\ \frac{S_n}{R_n}\right),\qquad \text{MAR} = 100\times\frac15\sum_n \text{NAR}_n .
\end{equation}
The cap at one stops a surplus of one nutrient from hiding a shortfall of another. MAR is reported in points out
of 100.

The requirements are:
\begin{itemize}
\item \textbf{Energy:} FAOSTAT's country- and year-specific Average Dietary Energy Requirement (a population
average).
\item \textbf{Protein and fat:} 10\% and 15\% of that energy (at 4 and 9\,kcal/g), the lower bounds of the WHO
goals \citep{who2003}.
\item \textbf{Zinc:} 11.9\,mg/day, the mean of the adult RNIs at low bioavailability (women 9.8, men 14.0;
\citealp{whofao2004} Table 12.4).
\item \textbf{Iron:} 21.55\,mg/day, the mean of the adult RNIs at 10\% bioavailability (women 29.4, men 13.7;
Table 13.5).
\end{itemize}
An alternative uses moderate-bioavailability zinc and 15\%-bioavailability iron.

Two cautions apply to reading these ratios.
\begin{itemize}
\item RNIs are upper-percentile \emph{individual} requirements, while $S_n$ is a national \emph{mean supply},
not intake. A ratio above one therefore does not mean that no one is deficient.
\item A ratio below one is a supply-side warning, not a measured deficiency.
\end{itemize}

\subsection{Standard yield shock and its bounds}\label{sec:shockmethod}
A yield change $\Delta Y<0$ (in \%) in crop $k$ removes part of the crop's contribution $C_{n,k}$ to national
supply. The fraction removed has two bounds.
\begin{itemize}
\item \textbf{Upper bound (no market response):} the loss passes through in proportion to self-sufficiency.
\item \textbf{Lower bound (exports absorb first):} exports absorb the lost production before domestic supply is
cut.
\end{itemize}
\begin{equation}\label{eq:shock}
f^{\text{up}} = \frac{|\Delta Y|}{100}\,\text{SSR}^{*},\qquad
f^{\text{low}} = \min\!\left(f^{\text{up}},\ \frac{\max\left(0,\ \tfrac{|\Delta Y|}{100}Q - X\right)}{\text{DS}}\right),\qquad
S'_n = S_n - f\,C_{n,k},
\end{equation}
where $Q$ is production, $X$ exports, DS domestic supply and $\text{SSR}=Q/\text{DS}$, all 2021--23 means. The
adjusted self-sufficiency $\text{SSR}^{*}$ is defined as follows:
\begin{itemize}
\item normally, $\text{SSR}^{*}=\min(1,\max(0,\text{SSR}))$;
\item when DS $\le0$ and $Q>0$, $\text{SSR}^{*}=1$. Such a balance occurs when exports, stock building or negative
residuals exceed the supply; a negative SSR would otherwise turn a loss into a gain.
\end{itemize}
The lower bound is never above the upper bound. Both rules were added in this audit (decision~E19;
Appendix~\ref{app:corrections}).

The change in each nutrient ($-100\,f\,C_{n,k}/S_n$, in \%) and the change in MAR are reported. The standard
shock is $\Delta Y=-10\%$ for every pair. It isolates the dietary channel, so countries and crops can be compared
on the same footing; it is not a projection.

\subsection{Historical pass-through}
For each pair, the percentage anomaly of the crop's food-energy supply is regressed on the same year's
percentage yield anomaly, both linearly detrended. This is done separately on the old FBS (1961--2013) and the
new FBS (2010--2023). The slope measures how much a bad harvest has historically reached food supply. Trade,
stocks, prices and income also move supply, so the slope is an association, not a causal effect.

\subsection{Food-supply diversity}\label{sec:diversitymethod}
The diversity of national food supply is measured, following \citet{remans2014}, by the effective number of
items supplying a nutrient:
\begin{equation}\label{eq:shannon}
H_n = -\sum_j p_{j,n}\ln p_{j,n},\qquad p_{j,n} = \frac{S_{n,j}}{\sum_{j'}S_{n,j'}},\qquad N^{\text{eff}}_n = e^{H_n},
\end{equation}
over the FBS food items $j$ (2021--2023 means). This describes national food \emph{supply}. It is not an
individual dietary diversity score, which needs individual or household intake data
(Section~\ref{sec:india}).

The share of a nutrient supplied by the nine study crops has:
\begin{itemize}
\item a \textbf{numerator}: the summed crop contributions (chickpea and lentil from SUA);
\item a \textbf{denominator}: the FBS national total.
\end{itemize}
It is reported only where SUA data exist; World has none, so its share is left blank.
